% ECONOMICS_PROBLEM: {"id": "C21-553", "title": "Long-run effects from proxies", "jel_code": "C21", "source_id": "OP-0429"}
% !TeX program = xelatex
\documentclass[11pt,a4paper]{article}
\usepackage[margin=23mm]{geometry}
\usepackage{fontspec}
\setmainfont{Latin Modern Roman}
\usepackage{amsmath,amssymb,mathtools}
\usepackage{xcolor,enumitem,booktabs,longtable,array,xurl,hyperref}
\definecolor{muted}{HTML}{53636C}
\hypersetup{colorlinks=true,urlcolor=blue,linkcolor=blue}
\setlength{\parindent}{0pt}
\setlength{\parskip}{5pt}
\newcommand{\R}{\mathbb R}
\newcommand{\E}{\mathbb E}
\newcommand{\N}{\mathbb N}
\newcommand{\ind}{\mathbf 1}
\newcommand{\Var}{\operatorname{Var}}
\newcommand{\Cov}{\operatorname{Cov}}
\newcommand{\supp}{\operatorname{supp}}
\DeclareMathOperator*{\argmin}{arg\,min}
\DeclareMathOperator*{\argmax}{arg\,max}
\newcommand{\entrymeta}[3]{{\sffamily\small\color{muted}\textbf{\texttt{#1}}\quad|\quad #2\par #3\par}}
\newcommand{\Problemref}[1]{\texttt{#1}}
\newcommand{\JELref}[2]{\texttt{#2} (JEL \texttt{#1})}
\begin{document}
\section*{Long-run effects from proxies}
\textbf{Problem ID:} \texttt{C21-553}\quad \textbf{JEL:} \texttt{C21}\par
% BEGIN ECONOMICS STATEMENT
\label{problem:OP-0429}
{\sffamily\small\color{muted}Original subject group: Econometrics, causal inference and measurement\par}
\label{definitions:problem:30007711}
\textbf{Audit classification:} Published research agenda. The2026 paper explicitly lists these broad challenges in Section6. The catalogue interval formula is editorial and already tractable under supplied tolerances; the earlier target-Q formulation omitted proxy-law transport from the new experiment to target users.
\subsection*{Definitions and precise research target}
Consider a new randomized platform experiment with treatment \(A\in\{0,1\}\), baseline covariates \(X\), short-run proxy \(Z\), and long-run potential outcome \(Y_H(a)\in[0,1]\) at horizon \(H\). Historical experiments observe \(A,X,Z,Y_H\); the new experiment initially observes only \(A,X,Z\). The target population has a specified covariate distribution \(Q\), possibly differing from historical participants, and the target is \(\tau_H=E_Q[Y_H(1)-Y_H(0)]\). Let \(m_a(x,z)\) be the historical conditional mean of \(Y_H\), and let \(\mathcal R_\delta\) restrict deviations between new and historical conditional means by \(\delta_a(x,z)\geq0\), including deviations from complete mediation and treatment-version changes. Assume treatment randomization, consistent measurement and overlap where historical extrapolation is used. Derive sharp bounds for \(\tau_H\) by integrating admissible means over each new treatment arm's proxy distribution and the target covariate distribution. Which long-run validation subsamples, persistent-treatment experiments or substantively credible restrictions justify smaller deviation bounds? Treat missing support separately with outcome bounds. The source explicitly identifies evolving treatments, incomplete proxy mediation and population mismatch as open challenges; this specified identified-set formulation adapts that agenda rather than assuming short-run surrogates universally identify persistent effects.

\textbf{Definitions, assumptions and mathematical qualifications.} Treatment \(a\) denotes a fully specified persistent intervention through horizon \(H\), including updates if any. Let \(K_a(dz\mid x)\) be the new-trial proxy law under randomized \(A=a\). To integrate it under target covariate law \(Q\), assume the target proxy law given \(X=x\) equals \(K_a\); otherwise supply bounds or target proxy data. Define historical predictive means \(m_a(x,z)=E[Y_H\mid A=a,X=x,Z=z,\text{historical}]\), not a causal effect of setting the proxy. Target conditional means obey \(l_a=\max(0,m_a-\delta_a)\leq\mu_a\leq u_a=\min(1,m_a+\delta_a)\). Off historical support set \(l_a=0,u_a=1\). With unrestricted conditional means within these bounds, the sharp effect interval is \([\int(\int l_1dK_1-\int u_0dK_0)dQ,\int(\int u_1dK_1-\int l_0dK_0)dQ]\). Computing this interval is elementary once kernels and tolerances are supplied; credible calibration and transport are the research challenges. Conditional kernel transport additionally requires target covariate absolute continuity with respect to the new experimental covariate law, and absolute continuity of target arm-specific \(X,Z\) laws with respect to historical arm-specific laws wherever \(m_a\) is used. Null-set failures use the full \([0,1]\) bounds rather than an arbitrary version of a conditional mean.
\subsection*{Variants and assumption changes}
\begin{enumerate}
\item \textbf{Editorial, stable fully mediated benchmark.} Assume a fixed treatment version, no direct long-run channel after conditioning on the proxy, stable conditional means and target proxy transport. Set \(\delta_a=0\) on common support.
\item \textbf{Editorial, evolving/persistent treatment.} Index outcomes by treatment paths \(Y_H(a_1,\ldots,a_H)\), allow direct long-run channels and population drift, and use representative long-run holdouts to calibrate deviations and proxy-law transport separately.
\end{enumerate}
\subsection*{References and source audit}
\textbf{Checked source.} arXiv:2608.08043v1 (8 August 2026), Section 6, propositions 6.3--6.6. Explicit research challenges concern evolving treatments, persistent effects, transport and decision accuracy. Full primary HTML and complete author list checked. \url{https://arxiv.org/html/2608.08043v1}.
\begin{enumerate}\raggedright
\item\hypertarget{definitions:source:30007711:1}{} Leif Sigerson; Tom Cunningham; Winston Chou; Sana Pandey; Jonathan Stray; Lo-Hua Yuan; Eytan Bakshy; Timothy Chan; Molly Davies; Maria Dimakopoulou; Simon Ejdemyr; Kenneth Hung; Nathan Kallus; Madhav Kumar; Thu Le; A. Demetri Pananos; Lee Richardson; Brennan Schaffner; Rose Tan; Martin Tingley; Nadia Tomova; Panagiotis Toulis; Wenjing Zheng; Zander Arnao; Dean Eckles. 2026. \emph{Evaluating for the Long Term: Learnings from Industry}. arXiv:2608.08043v1 (8 August 2026), Section 6, propositions 6.3--6.6.
\url{https://arxiv.org/html/2608.08043v1}.
DOI: \url{https://doi.org/10.48550/arXiv.2608.08043}.
\end{enumerate}

\subsection*{Common conventions: Source mathematical conventions}

These conventions apply to each entry unless explicitly replaced.

\textbf{Models and identification.} A primitive model \(m\) specifies all probability laws, feasible choices, information, equilibrium and measurement rules used by its target. The admissible class \(\mathcal M\) is fixed independently of outcomes. If \(P_O(m)\) is its observable probability law and \(\tau(m)\) the target, its identified set at observable law \(P\) is
\[
 \mathcal I_\tau(P)=\{\tau(m):m\in\mathcal M,\ P_O(m)=P\}.
\]
Point identification means that this set is a singleton. An empty set signals incompatibility of data and assumptions. Coordinate bounds need not describe the joint set. Bounds are sharp only if every retained target is attainable or arbitrarily approachable by compatible models. Finite bounds require explicit restrictions on unobserved quantities; unrestricted confounding, hidden wealth, extrapolation or arbitrary equilibrium selection can yield uninformative sets.

\textbf{Causal interventions.} A potential outcome \(Y_i(p)\) is the outcome under a complete policy regime \(p\), including timing, financing, information and market spillovers. The target population and units are fixed before treatment. With interference, use the complete assignment vector or a justified exposure mapping; individual no-interference notation alone cannot identify an economy-wide intervention. A difference of mean potential outcomes does not require the cross-world joint distribution. Conditioning on post-treatment migration, survival, enrollment or employment changes the estimand and needs separate justification.

\textbf{Statistical statements.} An estimator, confidence set, forecast or test specifies a sampling law, model class and loss/coverage criterion. Uniform \((1-\alpha)\)-coverage means \(\inf_{m\in\mathcal M}\Pr_m\{\tau(m)\in C_n\}\geq1-\alpha\), or an explicitly asymptotic version. The whole parameter space trivially has coverage, so informativeness requires a stated width, risk or power objective. A forecasting improvement is relative to a named benchmark, information set and evaluation population.

\textbf{Equilibrium and optimization.} An equilibrium comprises feasible plans, optimality, beliefs where needed, budget constraints and all market/resource clearing. The primitives or a stated benchmark define each correspondence \(\mathcal E(p;m)\). Nonemptiness is assumed or proved, never inferred just from compact policy sets. A selected-equilibrium target uses a declared selection rule; a robust target quantifies over all permitted equilibria. Use suprema or infima unless attainment is established. An \(\varepsilon\)-optimal policy is feasible and within \(\varepsilon\) of the finite optimum value. Report an empty feasible set separately.

\textbf{Welfare and accounting.} Normative utility, interpersonal normalization, social weights, numeraire, discounting and the use of public receipts are primitives. Behavioral utility need not coincide with the welfare criterion. Household welfare includes ownership income and rents once; adding profits again to compensating variation generally double-counts them. A monetary equivalent is defined by an explicit transfer and price convention, with monotonicity and range conditions. Average effects, mechanism contributions, transfers and welfare losses are different targets. Infinite sums require convergence and dynamic feasible plans require terminal or no-Ponzi conditions.

\textbf{Source versions.} A working-paper date, revision date and journal publication year are distinguished. A DOI for a journal version is not automatically the DOI of an earlier manuscript. Locators refer to the stated version and printed pages where specified. A metadata-only check does not verify a claim about a full-text conclusion. Every entry's source subsection narrows the provenance claim accordingly.
% END ECONOMICS STATEMENT
\end{document}
